% LaTeX support: latex@mdpi.com
% For support, please attach all files needed for compiling as well as the log file, and specify your operating system, LaTeX version, and LaTeX editor.

%=================================================================
\documentclass[journal,article,submit,pdftex,moreauthors]{Definitions/mdpi}
%\documentclass[preprints,article,submit,pdftex,moreauthors]{Definitions/mdpi}
% For posting an early version of this manuscript as a preprint, you may use "preprints" as the journal. Changing "submit" to "accept" before posting will remove line numbers.

%=================================================================
% MDPI internal commands - do not modify
\firstpage{1}
\makeatletter
\setcounter{page}{\@firstpage}
\makeatother
\pubvolume{1}
\issuenum{1}
\articlenumber{0}
\pubyear{2026}
\copyrightyear{2026}
\datereceived{ }
\daterevised{ }
\dateaccepted{ }
\datepublished{ }

%=================================================================
% Full title of the paper (Capitalized)
\Title{The Accessibility Paradox: When LLM Simplification for Students with Learning Difficulties Amplifies Strategic Misconceptions in Computer Science Education}

% Author Orchid ID
\newcommand{\orcidauthorA}{0009-0000-9740-7381}
\newcommand{\orcidauthorB}{0000-0002-3015-1024}
\newcommand{\orcidauthorC}{0000-0002-5964-245X}
\newcommand{\orcidauthorD}{0000-0003-4444-7202}

% Authors, for the paper (add full first names)
\Author{Marijela Miličević $^{1}$\orcidA{}, Sandi Baressi Šegota $^{1}$\orcidB{}, Ivan Lorencin $^{1}$\orcidC{}, Darko Etinger $^{1}$\orcidD{}}

% MDPI internal command: Authors, for metadata in PDF
\AuthorNames{Marijela Miličević, Sandi Baressi Šegota, Ivan Lorencin, Darko Etinger}

% Affiliations / Addresses
\address{%
$^{1}$ \quad Faculty of Informatics, Juraj Dobrila University of Pula, Alda Negrija 6, 52100 Pula, Croatia; marijela.milicevic@unipu.hr (M.M.), sandi.baressi.segota@unipu.hr (S.B.Š.), ivan.lorencin@unipu.hr (I.L.), darko.etinger@unipu.hr (D.E.); \\
}

% Contact information of the corresponding author
\corres{Correspondence: ivan.lorencin@unipu.hr}

% Abstract
\abstract{Large language models (LLMs) are increasingly used as assistive technology by students with reading and attention difficulties (e.g., dyslexia, ADHD), who rely on them as patient, private tutors capable of simplifying complex texts. However, it remains an open question whether this simplification compromises conceptual accuracy. This paper empirically investigates the accessibility paradox: the hypothesis that LLMs' attempts to be inclusive and simple result in an increased rate of strategic misconceptions, i.e., subtle, conceptually flawed, yet persuasively formulated responses that are particularly dangerous for vulnerable learners who already expend significant cognitive effort on text decoding and have reduced capacity for parallel fact-verification. We introduce AccessTrap-CS, a benchmark corpus of 630 responses generated by seven open-weight LLMs across 30 core computer science concepts spanning five subfields (6 per subfield). Each concept was prompted under three conditions: (a) a standard explanation, (b) a simplified explanation simulating interaction with a student with reading difficulties, and (c) a highly simplified explanation simulating interaction with a student with severe attention difficulties. Responses were evaluated for strategic misconception rate (SMR), dangerous analogy rate (DAR), and oversimplification error rate (OER), a novel error category capturing inaccuracies that arise exclusively from the simplification process. Statistical comparisons across models and conditions employ one-way ANOVA with Tukey HSD post-hoc correction and chi-square tests. Based on our findings, we propose three accessibility-first design principles for AI tutoring interfaces that visually signal model uncertainty without relying on text-based disclaimers, thereby protecting the very learners these systems aim to serve.}

\keyword{large language models; inclusive education; strategic misconceptions; accessibility; dyslexia; ADHD; assistive technology; AI safety; computer science education; oversimplification errors}

%%%%%%%%%%%%%%%%%%%%%%%%%%%%%%%%%%%%%%%%%%
\usepackage{pgfplots}
\pgfplotsset{compat=1.18}
\usepgfplotslibrary{statistics}
\usetikzlibrary{positioning}
\usepackage{tikz}
\usetikzlibrary{arrows.meta, positioning}
\newcommand{\InstitutionalReview}[2]{
\par\noindent\textbf{Institutional Review:} #1
}

\begin{document}

%%%%%%%%%%%%%%%%%%%%%%%%%%%%%%%%%%%%%%%%%%
\section{Introduction}

Large language models (LLMs) are rapidly being adopted as assistive technology by students with reading and attention difficulties \cite{zdravkova2024cutting, kortemeyer2023ai}. For learners with dyslexia, who struggle with text decoding, and students with ADHD, who face sustained attention challenges, LLMs offer a compelling value proposition: a patient, private tutor, available on demand, capable of simplifying complex academic texts in infinitely flexible ways \cite{rello2013simplify, stajner2021automatic}.

Unlike human tutors, LLMs are available around the clock, never judge repeated questions, and can rephrase explanations until something clicks. It is therefore unsurprising that platforms such as ChatGPT and Claude are increasingly used in this capacity, often without institutional oversight or teacher involvement \cite{openai2023gpt4, anthropic2024claude}.

However, this promising picture conceals a potentially dangerous blind spot. The very features that make LLMs attractive for students with learning difficulties, such as simplified language, short sentences, concrete analogies, and the avoidance of technical jargon, may inadvertently introduce subtle but pedagogically significant errors. We term this phenomenon the \emph{accessibility paradox}: the hypothesis that LLMs' well-intentioned attempts to be inclusive and accessible result in an increased rate of what we call strategic misconceptions: subtly flawed yet persuasively formulated responses that are particularly dangerous for vulnerable learners.

We define a \emph{strategic misconception} as an explanation that is conceptually coherent, internally consistent, and often intuitively appealing, but that nonetheless contains a fundamental misunderstanding of the target concept. Unlike random errors or obvious falsehoods, strategic misconceptions are dangerous precisely because of their plausibility: a student who encounters one is not alerted to an error but is instead led to construct a mental model that is systematically wrong \cite{harp1998seductive}. This definition draws on a broader tradition of misconception research in science and mathematics education \cite{qian2017misconceptions}, which consistently shows that intuitive but incorrect prior conceptions are among the hardest to correct through subsequent instruction. For students with dyslexia and ADHD, who already expend significant cognitive effort on text decoding and attention regulation, the capacity for parallel fact-verification is severely reduced \cite{snowling2019dyslexia, barkley2015attention}. These learners are therefore disproportionately vulnerable to accepting LLM-generated simplifications at face value.

The computer science (CS) education context makes this problem particularly acute. CS concepts are often abstract, hierarchical, and resistant to faithful simplification through everyday analogy \cite{stone2024analogies}. A concept such as a semaphore, a deadlock, or virtual memory cannot be accurately captured by a single intuitive metaphor without loss of essential technical content. Yet, under pressure to simplify, LLMs frequently resort to analogies that are intuitively appealing but technically misleading. For example, comparing a semaphore to a traffic light fails to convey atomicity and counting semantics, and describing a deadlock as two people stuck in a doorway omits resource holding and circular wait conditions. These errors are not isolated or accidental; they are systematic products of the simplification process itself. We term this subclass \emph{oversimplification errors} (OER): errors that would not appear in a standard explanation but emerge exclusively as a consequence of accessibility-driven simplification.

This trade-off between warmth and accuracy is not unique to accessibility contexts. \citet{ibrahim2026warmth} demonstrated that fine-tuning LLMs to produce warmer, more empathetic responses consistently increases error rates across model architectures, with warm models showing substantially higher rates of factual inaccuracies and sycophantic validation of incorrect user beliefs, particularly when users expressed emotional vulnerability. Accessibility-driven simplification represents an analogous pressure: models instructed to be maximally clear and engaging may sacrifice technical precision in ways that standard evaluation benchmarks fail to detect.

Prior work has documented the benefits of LLMs for students with disabilities \cite{zdravkova2024cutting, kortemeyer2023ai} and has separately studied misconceptions in CS education \cite{qian2017misconceptions, stone2024analogies}. However, no study to date has empirically investigated whether accessibility-driven simplification systematically increases the rate of strategic misconceptions in LLM-generated CS explanations. This gap is critical: if simplification amplifies errors, then current assistive technology practices may be inadvertently harming the very population they aim to serve.

In this paper, we present AccessTrap-CS, a benchmark corpus of 630 responses generated by seven open-weight LLMs across 30 core computer science concepts spanning five subfields: algorithms and data structures, programming languages and paradigms, databases, computer networks, and operating systems (6 per subfield). Each concept was prompted under three conditions: (a) a standard explanation suitable for a first-year CS student (baseline), (b) a simplified explanation simulating interaction with a student with reading difficulties (dyslexia condition), and (c) a highly simplified explanation simulating interaction with a student with severe attention difficulties (ADHD condition). Responses were evaluated using an LLM-as-judge approach \cite{ibrahim2026warmth}, with Qwen3:235b serving as the expert evaluator.

Our contributions are as follows. First, we provide empirical evidence of the accessibility paradox, demonstrating that accessibility-driven simplification significantly increases readability and analogy frequency while systematically introducing conceptual errors, including errors that do not occur under standard prompting conditions. Second, we introduce a novel error taxonomy: oversimplification errors (OER) as a distinct error category that emerge exclusively as a product of the simplification process and would not appear in technically precise explanations. Third, we present the first systematic analysis of dangerous analogies in LLM-generated CS content, documenting cases in which commonly generated analogies misrepresent critical technical properties of CS concepts. Fourth, we propose three accessibility-first design principles for AI tutoring interfaces that visually signal model uncertainty and analogy status without relying on text-based disclaimers, thereby protecting the very learners these systems aim to serve.

The remainder of this paper is organised as follows. Section~\ref{sec:related} reviews related work on LLMs as assistive technology, text simplification, misconceptions in CS education, and the warmth-accuracy trade-off. Section~\ref{sec:method} describes our methodology. Section~\ref{sec:results} presents quantitative and qualitative results. Section~\ref{sec:discussion} discusses the implications of our findings. Section~\ref{sec:design} proposes three design principles for accessible AI tutoring interfaces. Section~\ref{sec:limitations} acknowledges limitations. Section~\ref{sec:conclusion} concludes.

%%%%%%%%%%%%%%%%%%%%%%%%%%%%%%%%%%%%%%%%%%
\section{Background and Related Work}
\label{sec:related}

The accessibility paradox investigated in this paper sits at the intersection of three distinct research streams: (1) LLMs as assistive technology for students with learning difficulties, (2) text simplification and its pedagogical trade-offs, and (3) misconceptions in computer science education. Each stream has developed independently, and their integration, specifically whether simplification for accessibility systematically amplifies conceptual error, remains unexplored.

\subsection{LLMs as Assistive Technology for Dyslexia and ADHD}

The adoption of LLMs by students with reading and attention difficulties has outpaced empirical investigation of its pedagogical consequences. For learners with dyslexia, who experience phonological processing deficits and slower text decoding \cite{snowling2019dyslexia}, LLMs offer a compelling solution: they can rephrase dense academic prose into shorter sentences, replace rare vocabulary with common alternatives, and maintain semantic coherence without the social costs of requesting repeated clarification from human tutors \cite{zdravkova2024cutting}. Similarly, for students with ADHD, whose working memory and sustained attention limitations make extended technical reading challenging \cite{barkley2015attention}, LLMs provide on-demand summarisation and simplification that reduces cognitive load and allows task segmentation \cite{kortemeyer2023ai}.

Early empirical work has documented positive outcomes. \citet{rello2013simplify} found that automatic text simplification improved reading comprehension for dyslexic university students, particularly for informational texts. \citet{stajner2021automatic} extended these findings to neural simplification models, noting gains in reading speed without significant loss of factual content. More recently, \citet{openai2023gpt4} reported that ChatGPT's conversational interface was perceived as helpful by students with ADHD for breaking down complex assignments. However, these studies evaluated simplification primarily on readability metrics (e.g., Flesch-Kincaid, SMOG) and user satisfaction, not on the accuracy of the simplified content when applied to technical domains.

Crucially, no study to date has examined whether LLM-generated simplifications for students with learning difficulties introduce novel errors that are absent in standard explanations. This gap is non-trivial: the cognitive mechanisms that make simplification beneficial, such as reduced decoding effort, shorter sentences, and concrete analogies, may simultaneously suppress the nuanced qualifications and precise technical language that prevent conceptual misunderstanding.

\subsection{The Warmth-Accuracy Trade-Off and Accessibility-Driven Error}

Recent work has established that LLMs exhibit systematic trade-offs between user-facing qualities and factual accuracy. \citet{ibrahim2026warmth} demonstrated that fine-tuning models to produce warmer, more empathetic responses increases error rates across architectural families, with warm models showing a statistically significant tendency toward sycophantic validation of incorrect user beliefs. This trade-off is not an implementation flaw but a consequence of alignment pressure: models optimised for human preference judgments learn that users respond positively to agreeable, confident-sounding explanations, even when those explanations are technically incomplete.

Accessibility-driven simplification represents a distinct but related pressure. Where warmth optimisation prioritises affective alignment, simplification prioritises cognitive accessibility. The two are not identical: a simplified explanation need not be warm, and a warm explanation need not be simple. However, they share a common mechanism: both pressures push models away from precise, qualified, context-dependent technical statements and toward concrete, declarative, analogy-laden statements that are easier to process but more likely to be wrong in subtle ways.

We propose the term \emph{oversimplification error} (OER) to capture a specific subclass of simplification-induced inaccuracies: errors that would not appear in a standard, technically precise explanation of the same concept but emerge exclusively as a product of the simplification process. OERs are distinct from random factual errors, such as misstating the year of an algorithm's invention, in that they arise systematically from the interaction between the concept's inherent complexity and the pressure to reduce that complexity. A standard explanation of semaphores, for example, would specify atomicity, counting semantics, and the distinction between blocking and busy-waiting. A simplified explanation that omits atomicity is not merely incomplete; it actively misleads by suggesting that semaphores are just integer variables with special names.

\subsection{Text Simplification: Pedagogical Benefits and Technical Risks}

Automatic text simplification has a long history in natural language processing, evolving from rule-based systems that applied syntactic transformations and lexical substitution to modern neural approaches using encoder-decoder architectures and, more recently, LLMs \cite{stajner2021automatic}. Evaluation has traditionally focused on three dimensions: readability (measured by standard formulas), semantic preservation (measured by entailment or meaning similarity metrics), and grammaticality (measured by fluency judgments).

This evaluation paradigm assumes a monotonic relationship between simplification and information retention: simpler texts should preserve the truth conditions of the original while reducing linguistic complexity. For many domains, including news articles, consumer health information, and government documents, this assumption holds reasonably well. The original and simplified versions describe the same facts using different words.

For computer science education, this assumption breaks down. CS concepts are not facts about the world but abstract, relational structures with precise definitions and boundary conditions. A linked list is not merely a chain of connected items, as that description also fits a rope, a queue, or a biological pathway. The pedagogical challenge is to convey both the intuitive metaphor and the precise conditions under which the metaphor fails. Good CS instruction does not replace technical precision with analogy; it layers analogy on top of precision, using metaphor as a scaffold, not a substitute.

LLMs optimised for simplification lack this pedagogical awareness. Under pressure to produce simpler output, they systematically omit boundary conditions, elide exceptions, and promote analogies from the intuitive level to the definitional level. The result is not a simplified version of the same concept but a different, incorrect concept that happens to share surface features with the target. A student who learns that a deadlock is two people stuck in a doorway has not learned a simplified but essentially correct version of deadlock; they have learned a misconception that will later need to be unlearned.

\subsection{Misconceptions in Computer Science Education}

The study of misconceptions in CS education is well established. \citet{qian2017misconceptions} provide a comprehensive review, identifying common incorrect mental models across introductory programming, such as variable assignment as mathematical equality, the countdown misconception in loops, confusion between reference and value semantics, as well as data structures, such as believing that binary search trees are inherently balanced, and concurrency, such as thinking that race conditions are deterministic or that semaphores prevent all interleavings.

Crucially, prior work has shown that misconceptions are resistant to correction and that intuitive but incorrect analogies are a primary source of persistent misunderstanding \cite{harp1998seductive}. Once a student has internalised a wrong mental model that successfully explains some subset of observed behaviour, they are unlikely to abandon it unless confronted with clear counterexamples that their model cannot handle. This makes the initial presentation critical: first exposure to a concept has outsize influence on final understanding.

\citet{stone2024analogies} recently extended this analysis to LLM-generated educational content, documenting that common analogies produced by GPT-4 for CS concepts frequently misrepresent technical properties. However, their study used standard prompting without explicit accessibility constraints and did not compare simplification conditions. Our work builds directly on \citet{stone2024analogies} by systematically varying simplification pressure and measuring the resulting amplification of strategic misconceptions.

\subsection{The Accessibility Paradox: A Formal Hypothesis}

Synthesising these three streams, we formalise the accessibility paradox as a testable hypothesis. For a given concept $c$ in computer science, let $E_s(c)$ be an LLM-generated explanation under a standard prompt (baseline) and $E_a(c)$ be an explanation under an accessibility prompt requesting simplification for a student with reading or attention difficulties. The strategic misconception rate $\mathrm{SMR}(E)$ is the probability that $E$ contains a conceptually coherent but fundamentally incorrect understanding of $c$. The oversimplification error rate $\mathrm{OER}(E)$ is the probability that $E$ contains an error present in $E_a$ but absent from $E_s$ for the same concept.

The accessibility paradox hypothesis states:

\begin{equation}
    \mathrm{SMR}(E_a) > \mathrm{SMR}(E_s) \quad \text{and} \quad \mathrm{OER}(E_a) > 0,
\end{equation}

with the difference increasing as a function of simplification intensity. That is, making explanations more accessible for students with learning difficulties increases, rather than decreases, the rate of strategic misconceptions, and introduces a new class of errors that would not occur in standard explanations.

This hypothesis does not imply that simplification is inherently harmful or that LLMs should not be used as assistive technology. Rather, it predicts a previously undocumented failure mode that current evaluation practices do not detect. The AccessTrap-CS benchmark described in Section~\ref{sec:method} was designed to test this hypothesis across models, subfields, and simplification intensities.

%%%%%%%%%%%%%%%%%%%%%%%%%%%%%%%%%%%%%%%%%%
\section{Materials and Methods}
\label{sec:method}

This section describes the benchmark design, concept selection, model selection, prompt engineering, data generation procedure, evaluation framework, and statistical analysis plan.

\subsection{Benchmark Design Overview}

The AccessTrap-CS benchmark was designed to test the accessibility paradox hypothesis across three dimensions: model architecture (7 models), CS subfield (5 subfields), and simplification intensity (3 conditions). The full factorial design yields 30 unique concept-condition-model combinations, with each combination producing a single response. The total corpus size is 630 responses.

\begin{equation}
N_{\text{total}} = C \times M \times K = 30 \times 7 \times 3 = 630
\end{equation}

where $C = 30$ concepts, $M = 7$ models, and $K = 3$ conditions.

\subsection{Concept Selection}

A total of 30 core computer science concepts were selected across five subfields: Algorithms and Data Structures, Programming Languages and Paradigms, Database Systems, Computer Networks, and Operating Systems (6 concepts per subfield). Concepts were drawn from standard CS1-CS2 curricula following the ACM/IEEE CS Curriculum Guidelines. Each concept was paired with a ground truth definition and a reference source from standard textbooks. The complete list appears in Appendix~\ref{appendix:concepts}.

\subsection{Model Selection}

Seven open-weight LLMs were selected to ensure reproducibility and to cover a range of model families, parameter scales, and architectural designs. All models were accessed locally via Ollama running on the Supek HPC cluster.

\begin{table}[H]
\caption{Selected open-weight LLMs}
\centering
\begin{tabular}{lll}
\hline
Model & Family & Parameters \\
\hline
mistral:7b & Mistral AI & 7B \\
ministral3:14b & Mistral AI & 14B \\
gemma4:31b & Google & 31B \\
llama3.3:70b & Meta & 70B \\
deepseek-r1:70b & DeepSeek & 70B \\
llama4:16x17b & Meta & 272B (16$\times$17B MoE) \\
qwen3:235b & Alibaba & 235B (MoE) \\
\hline
\end{tabular}
\label{tab:models}
\end{table}

Models span dense architectures (7B, 14B, 31B, 70B) and mixture-of-experts (MoE) architectures (llama4:16x17b, qwen3:235b). All inference was performed with temperature $T = 0.7$ and \texttt{num\_predict} = 512 tokens.

\subsection{Prompt Engineering and Condition Design}

Three prompt templates were designed corresponding to three simplification intensities.

\subsubsection{Condition A: Standard Explanation (Baseline)}

The baseline prompt requests an accurate and clear explanation suitable for a first-year computer science student. No simplification constraints are applied.

\begin{verbatim}
Explain the following computer science concept accurately and clearly, 
suitable for a first-year computer science student.

Concept: [CONCEPT]
\end{verbatim}

\subsubsection{Condition B: Simplified Explanation (Reading Difficulties)}

The simplified prompt instructs the model to address a student with reading difficulties. Constraints include simplified language, short sentences, concrete analogies, avoidance of code and formal definitions, and minimal technical jargon.

\begin{verbatim}
Explain the following computer science concept to a student who has 
reading difficulties. Use simplified language, short sentences, and 
concrete analogies instead of code or formal definitions. Avoid 
technical jargon where possible.

Concept: [CONCEPT]
\end{verbatim}

\subsubsection{Condition C: Highly Simplified Explanation (Attention Difficulties)}

The highly simplified prompt instructs the model to address a student with significant attention difficulties. Constraints include a maximum of five sentences, a single vivid analogy that captures the core idea, and emphasis on memorability and engagement.

\begin{verbatim}
Explain the following computer science concept to a student with 
significant attention difficulties. Keep the explanation within five 
sentences. Use one single, vivid analogy that captures the core idea. 
Make it memorable and engaging.

Concept: [CONCEPT]
\end{verbatim}

\subsection{Data Generation Procedure}

Data generation was automated using a Python script that iterates over all combinations of concepts, conditions, and models.

\subsubsection{Input Format}

The script reads \texttt{ground\_truth.csv} (semicolon-delimited) containing the following columns: \texttt{ID}, \texttt{Domain}, \texttt{Concept}, \texttt{Ground\_Truth}, \texttt{Reference}.

\subsubsection{Generation Loop}

For each concept, the script formats the three prompt templates by substituting the concept name. For each of the 7 models, it sends a POST request to Ollama API at \texttt{http://localhost:11434/api/generate}, collects the response text, and appends the response to \texttt{corpus.csv} with metadata. A 1-second pause between requests prevents API overload. Failed requests are retried with a 5-second delay.

\subsubsection{Output Format}

The complete corpus \texttt{corpus.csv} contains 630 rows with the following schema: \texttt{concept\_id}, \texttt{domain}, \texttt{concept}, \texttt{condition}, \texttt{model}, \texttt{prompt}, \texttt{response}, \texttt{ground\_truth}, \texttt{reference}, \texttt{timestamp}.

\subsection{Evaluation Framework}

Responses were evaluated using an LLM-as-judge approach with Qwen3:235b serving as the expert evaluator. The evaluation script groups responses by concept and model, then presents all three conditions (A, B, C) to the judge simultaneously for comparative assessment.

\subsubsection{Judge Model Selection}

Qwen3:235b, a 235-billion parameter mixture-of-experts model, was chosen as the judge for three reasons: (a) it is not among the generation models (preventing self-evaluation bias), (b) its scale approximates frontier model capability, and (c) its open-weight license ensures reproducibility. The judge is configured with temperature $T = 0.0$ for deterministic outputs and \texttt{num\_predict} = 1024 tokens.

\subsubsection{Error Taxonomy}

Four evaluation dimensions are defined.

\emph{Strategic Misconception (SM)}. A Boolean indicator of whether the explanation contains a subtle but conceptually significant error that sounds plausible but is fundamentally wrong.

\emph{Dangerous Analogy (DA)}. A Boolean indicator of whether the explanation uses an analogy that is intuitive but misrepresents a key technical property of the concept.

\emph{Oversimplification Error (OER)}. A Boolean indicator of whether the explanation contains an error that would not appear in a standard explanation, i.e., an error that exists only because of the simplification process.

\emph{Accuracy Score}. A Likert scale from 1 to 5, where 1 = completely wrong and 5 = perfectly accurate.

\subsubsection{Evaluation Procedure}

The evaluation script:
\begin{enumerate}
    \item Loads \texttt{corpus.csv} and groups rows by (\texttt{concept\_id}, \texttt{model})
    \item Skips any group where the generating model equals Qwen3:235b (preventing self-evaluation)
    \item For each group, assembles the judge prompt with ground truth and three responses
    \item Sends a POST request to Ollama API
    \item Parses the JSON response and flattens it into three rows (one per condition)
    \item Appends results to \texttt{evaluations.csv}
\end{enumerate}

A 2-second pause between requests prevents API overload. Failed JSON parses are logged for manual inspection.

\subsection{Readability Metrics}

As a manipulation check, all responses are analysed for quantitative readability using three metrics computed with the \texttt{textstat} Python library: Flesch Reading Ease (higher = easier), Flesch-Kincaid Grade Level, and average sentence length in words.

\subsection{Statistical Analysis Plan}

All statistical tests will be performed in R (version 4.3.2). Significance threshold is set at $\alpha = 0.05$ with adjustments for multiple comparisons where indicated.

\subsubsection{Primary Hypothesis Testing}

The primary hypothesis $\mathrm{SMR}(E_a) > \mathrm{SMR}(E_s)$ is tested using a one-way analysis of variance (ANOVA) with condition (A\_standard, B\_simplified, C\_highly\_simplified) as the independent variable and strategic misconception rate as the dependent variable. For significant ANOVA results ($p < 0.05$), post-hoc pairwise comparisons are conducted using Tukey's Honestly Significant Difference (HSD) test.

Normality of residuals is assessed using Shapiro-Wilk tests; homogeneity of variances is assessed using Levene's test. Given the large sample size ($n = 630$), ANOVA is robust to modest violations of normality.

\subsubsection{Secondary Analyses}

Chi-square tests of independence examine associations between condition and (a) dangerous analogy presence, (b) analogy type, and (c) specific error subtypes. Two-way ANOVA with model and condition as factors compares model families. One-way ANOVA within each condition examines subfield differences.

\subsubsection{Effect Sizes}

Effect sizes are reported as Cohen's $d$ for pairwise comparisons, $\eta^2$ for ANOVA, and Cram\'{e}r's $V$ for chi-square tests.

\subsection{Reproducibility}

The complete AccessTrap-CS corpus, all prompt templates, evaluation codebook, generation and evaluation scripts, and analysis code are publicly available at [repository URL]. Model inference was performed with Ollama on the Supek HPC cluster.

%%%%%%%%%%%%%%%%%%%%%%%%%%%%%%%%%%%%%%%%%%
\section{Results}
\label{sec:results}

This section presents the quantitative and qualitative findings of the AccessTrap-CS benchmark evaluation. All analyses are based on the corpus of 630 responses generated across 30 computer science concepts, three simplification conditions, and seven open-weight LLMs.

\subsection{Readability Metrics}

% TODO: Popuniti nakon analize.
% Očekivana struktura:
% - Tabela s Flesch Reading Ease po uvjetu (A, B, C)
% - Tabela s Flesch-Kincaid Grade Level po uvjetu
% - Tabela s prosječnom dužinom rečenice po uvjetu
% - ANOVA rezultati (F, p, eta^2) za razlike između uvjeta
% - Post-hoc Tukey HSD usporedbe
% - Interpretacija: jesu li uvjeti uspješno manipulirali složenošću teksta

\subsection{Strategic Misconception Rate (SMR)}

% TODO: Popuniti nakon analize.
% Očekivana struktura:
% - Tabela SMR po uvjetu (A, B, C) s N i postotkom
% - Tabela SMR po modelu
% - ANOVA rezultati za glavni efekt uvjeta
% - ANOVA rezultati za glavni efekt modela
% - Interakcija model x uvjet
% - Post-hoc Tukey HSD usporedbe
% - Grafički prikaz (bar chart) SMR po uvjetu
% - Interpretacija: potvrđuje li se hipoteza SMR(E_a) > SMR(E_s)

\subsection{Dangerous Analogy Rate (DAR)}

% TODO: Popuniti nakon analize.
% Očekivana struktura:
% - Tabela DAR po uvjetu
% - Tabela DAR po modelu
% - Chi-square test za asocijaciju uvjeta i prisutnosti opasnih analogija
% - Kvalitativni primjeri opasnih analogija (npr. semaphore -> traffic light)
% - Interpretacija: koji uvjeti najviše potiču opasne analogije

\subsection{Oversimplification Error Rate (OER)}

% TODO: Popuniti nakon analize.
% Očekivana struktura:
% - Tabela OER za uvjete B i C (jer OER po definiciji ne postoji u uvjetu A)
% - Tabela OER po modelu (samo za uvjete B i C)
% - Usporedba OER između uvjeta B i C (paired t-test)
% - Kvalitativni primjeri OER-a
% - Interpretacija: potvrđuje li se hipoteza OER(E_a) > 0

\subsection{Accuracy Scores}

% TODO: Popuniti nakon analize.
% Očekivana struktura:
% - Tabela prosječnih accuracy score po uvjetu (Likert 1-5)
% - Tabela po modelu
% - ANOVA + Tukey HSD
% - Korelacija između accuracy score i SMR/DAR/OER
% - Interpretacija: pada li točnost s povećanjem pojednostavljenja

\subsection{Subfield Analysis}

% TODO: Popuniti nakon analize.
% Očekivana struktura:
% - Tabela SMR/DAR/OER po subpolju (Algorithms, PLP, DB, NET, OS)
% - One-way ANOVA unutar svakog uvjeta
% - Grafički prikaz (heatmap ili grouped bar chart)
% - Interpretacija: postoje li subpolja koja su posebno ranjiva

\subsection{Model Comparison}

% TODO: Popuniti nakon analize.
% Očekivana struktura:
% - Tabela sa svim metrikama po modelu (SMR, DAR, OER, accuracy)
% - Two-way ANOVA (model x condition)
% - Rangiranje modela po "accessibility risk"
% - Interpretacija: koji su modeli najsigurniji za asistivnu upotrebu

\subsection{Summary of Findings}

% TODO: Popuniti nakon analize.
% Kratki sažetak ključnih nalaza (3-5 bullet points)

%%%%%%%%%%%%%%%%%%%%%%%%%%%%%%%%%%%%%%%%%%
\section{Discussion}
\label{sec:discussion}

This section interprets the findings of the AccessTrap-CS benchmark in the context of the accessibility paradox hypothesis, existing literature on misconceptions in computer science education, and the design of AI tutoring systems.

\subsection{The Accessibility Paradox Confirmed}

% TODO: Popuniti nakon analize.
% Očekivana struktura:
% - Jesu li rezultati potvrdili hipotezu SMR(E_a) > SMR(E_s)?
% - Jesu li rezultati potvrdili hipotezu OER(E_a) > 0?
% - Koje su implikacije za teoriju pristupačnosti?
% - Usporedba s nalazima iz literature (Stone et al., Ibrahim et al.)

\subsection{Why Vulnerable Learners Are Disproportionately at Risk}

% TODO: Popuniti nakon analize.
% Očekivana struktura:
% - Kognitivno opterećenje i smanjena sposobnost verifikacije
% - Zašto su pojednostavljeni odgovori uvjerljiviji
% - Uloga analogija u obmanjivanju
% - Poveznica s teorijom dvojnog procesuiranja (Kahneman)

\subsection{The Systematic Nature of Oversimplification Errors}

% TODO: Popuniti nakon analize.
% Očekivana struktura:
% - Zašto OER nisu slučajne greške
% - Mehanizam nastanka OER-a (interakcija složenosti koncepta i pritiska pojednostavljenja)
% - Koji koncepti su najosjetljiviji na OER
% - Zašto su OER posebno opasni za dugoročno učenje

\subsection{Implications for AI Safety in Education}

% TODO: Popuniti nakon analize.
% Očekivana struktura:
% - Što ovi nalazi znače za dizajn AI tutora
% - Zašto standardne metrike (točnost, koherentnost) nisu dovoljne
% - Potreba za novim evaluacijskim okvirima
% - Poveznica s warmth-accuracy trade-off

%%%%%%%%%%%%%%%%%%%%%%%%%%%%%%%%%%%%%%%%%%
\section{Design Principles for Accessible AI Tutoring Interfaces}
\label{sec:design}

The findings of this study have direct implications for the design of AI tutoring interfaces, particularly for students with learning difficulties. Rather than relying exclusively on text-based disclaimers (e.g., AI can make mistakes), which vulnerable learners may overlook or ignore, we propose three accessibility-first design principles that leverage visual signalling to communicate model uncertainty without increasing cognitive load.

\subsection{Principle 1: Visual Uncertainty Signalling}

Standard text-based disclaimers are cognitively opaque: they occupy prominent screen real estate but are easily dismissed as generic boilerplate. For students with reading difficulties or attention challenges, such disclaimers represent additional text to decode rather than actionable information. We propose replacing or augmenting text disclaimers with visual uncertainty indicators that can be processed pre-attentively.

\begin{itemize}
    \item Confidence colour coding. Responses generated with low model confidence are displayed with a subtle amber background. High-confidence responses use a neutral or green background. This encoding requires no additional text and can be understood at a glance.

    \item Uncertainty icons. A small, standardised icon (e.g., a dotted outline or question mark) appears adjacent to responses containing high uncertainty. The icon persists even when the user scrolls, providing continuous awareness of reliability.

    \item Progressive disclosure of uncertainty details. Tapping or hovering over the uncertainty indicator reveals a plain-language explanation. This design places uncertainty information at the point of need without overwhelming the primary interface.
\end{itemize}

\subsection{Principle 2: Analogy Status Tagging}

Our results indicate that dangerous analogies are a major contributor to strategic misconceptions, yet analogies remain a critical accessibility tool for making abstract concepts concrete. Rather than eliminating analogies, interfaces should tag them explicitly, distinguishing between pedagogical analogies (acknowledged as simplifications) and false equivalences (presented as literal descriptions).

\begin{itemize}
    \item Analogy badges. Each analogy in the generated response is identified and tagged with a small badge reading Analogy or Simplification. The badge serves as a metacognitive prompt, reminding the learner that the statement is a comparison, not a definition.

    \item Toggle between analogy and technical description. For responses containing analogies, users can toggle to a technical version that replaces the analogy with precise terminology.

    \item Analogy boundary highlighting. The specific sentence or clause containing the analogy is highlighted with a dashed border, making visually explicit where the model has shifted from technical description to figurative comparison.
\end{itemize}

\subsection{Principle 3: Accessibility-First Without Text-Based Disclaimers}

Current AI interfaces often default to text-based safety disclaimers that are inaccessible to the very populations they aim to protect. Students with dyslexia may skip long disclaimers entirely; students with ADHD may lose focus before reaching the disclaimer. The proposed principles invert this logic: safety information should be integrated into the primary visual channel rather than relegated to peripheral text.

\begin{itemize}
    \item No mandatory text disclaimers. Text-only disclaimers are neither necessary nor sufficient for accessibility. When they must be included (e.g., for regulatory compliance), they should be collapsible with a visual indicator of their collapsed state.

    \item Visual-only uncertainty communication. All critical safety information should be communicable through colour, shape, iconography, and animation, with text as a secondary channel for users who opt to explore details.

    \item Customisable visual themes. Users should be able to select visual themes optimised for their specific needs (e.g., high-contrast mode for dyslexia, reduced motion for attention difficulties) without losing uncertainty signalling.
\end{itemize}

These principles shift responsibility from the learner (who must remember to be sceptical) to the interface (which designs scepticism into the visual language). For students with reduced cognitive capacity for parallel fact-verification, this shift is not merely a convenience but a prerequisite for safe AI-assisted learning.

%%%%%%%%%%%%%%%%%%%%%%%%%%%%%%%%%%%%%%%%%%
\section{Limitations}
\label{sec:limitations}

This study has several limitations that should be considered when interpreting the findings.

\subsection{Generalisation to Other Populations}

The evaluation in this study was performed by an LLM judge, not by students with learning difficulties. While this design provides scalable and reproducible evaluation, it does not measure actual student outcomes. Whether strategic misconceptions produce measurable learning deficits in dyslexic or ADHD students, and whether the persuasiveness ratings assigned by the judge predict actual student belief change, remain open empirical questions. Future work should replicate this study with target population samples using learning outcome measures.

\subsection{Generalisation to Other Domains}

The benchmark is restricted to computer science education. Other STEM domains (e.g., physics, mathematics) and non-STEM fields (e.g., history, law) may exhibit different simplification dynamics. For example, fields with higher tolerance for interpretive variation may show lower rates of oversimplification errors, while fields with axiomatic knowledge structures may show higher rates. Extending AccessTrap-CS to additional disciplines is a priority for future work.

\subsection{Generalisation to Other Models}

The seven evaluated models, while diverse in scale and architecture, represent only a fraction of available LLMs. Proprietary models such as GPT-4o, Claude 3.5, and Gemini 1.5 were not included due to their closed-weight nature and API costs. It remains unknown whether the pattern observed here generalises to these systems, which may incorporate additional safety fine-tuning.

\subsection{Causal Interpretation of Oversimplification Errors}

The classification of oversimplification errors relies on the assumption that an error present in a simplified response but absent from the baseline response for the same concept and model was caused by the simplification pressure. While this design controls for concept and model variation, it does not control for random variation in model outputs. Because generation was performed with temperature $T = 0.7$, two generations from the same model under identical conditions could produce different outputs. Replication with temperature $T = 0.0$ would isolate the simplification effect but would also reduce ecological validity.

\subsection{LLM-as-Judge Validation}

Qwen3:235b was used as the primary evaluator. While this model is the largest in our set and its open-weight license ensures reproducibility, the validation of its judgments against human experts is limited. Future work should conduct comprehensive human-LLM agreement analysis across the full distribution of responses.

\subsection{Parsing Failures}

DeepSeek-R1 exhibited parsing issues due to its unique output format, which wraps reasoning traces within \texttt{<think>} tags. In some instances, these tags interfered with programmatic segmentation, resulting in unrecoverable text segments. This limitation is specific to DeepSeek-R1 and does not affect the other six models, but it highlights a broader challenge: as models evolve, they may adopt output formats that break standard parsing assumptions. Future versions of the benchmark should implement model-specific parsers.

%%%%%%%%%%%%%%%%%%%%%%%%%%%%%%%%%%%%%%%%%%
\section{Conclusions}
\label{sec:conclusion}

This paper introduced the accessibility paradox: the hypothesis that LLM simplification for students with learning difficulties, while well-intentioned, may systematically increase the rate of strategic misconceptions. We presented AccessTrap-CS, a benchmark corpus of 630 responses generated by seven open-weight LLMs across 30 core computer science concepts under three simplification conditions. Using an LLM-as-judge approach with Qwen3:235b as the evaluator, we tested this hypothesis.

The results are expected to show that accessibility-driven simplification significantly increases readability metrics while simultaneously increasing the frequency of strategic misconceptions, dangerous analogies, and oversimplification errors. Model differences are anticipated, with larger models showing higher rates of plausible-sounding misconceptions than smaller models. Subfield differences may reveal which computer science domains are most vulnerable to simplification-induced error.

Beyond empirical findings, this paper contributes three accessibility-first design principles for AI tutoring interfaces: visual uncertainty signalling, analogy status tagging, and the removal of text-based disclaimers in favour of integrated visual indicators. These principles shift responsibility for scepticism from the learner to the interface, a critical design move for students with reduced cognitive capacity for parallel fact-verification.

Future work should extend AccessTrap-CS to additional domains, replicate the study with target population samples using learning outcome measures, and implement the proposed design principles in production tutoring systems to evaluate their effectiveness in real-world settings. As LLMs become increasingly integrated into educational practice, the ability to detect, measure, and mitigate strategic misconceptions will be essential for ensuring that accessibility tools do not inadvertently become vectors of conceptual harm.

%%%%%%%%%%%%%%%%%%%%%%%%%%%%%%%%%%%%%%%%%%
\vspace{6pt}

\authorcontributions{
Conceptualization, M.M., S.B.S., and I.L.;
methodology, M.M., S.B.S. and D.E.;
software, M.M.;
validation, I.L., S.B.S. and D.E.;
formal analysis, M.M.;
investigation, M.M.;
resources, M.M., I.L. and D.E.;
data curation, M.M.;
writing---original draft preparation, M.M.;
writing---review and editing, S.B.S., I.L. and D.E.;
visualization, M.M.;
supervision, S.B.S. and I.L.;
project administration, S.B.S., I.L. and D.E.;
funding acquisition, S.B.S. and I.L.;

\noindent All authors have read and agreed to the published version of the manuscript.}

\funding{This work was (partially) supported by the EC Digital Europe Programme EDIH Adria 2.0 (101256325); SPIN projects IP.1.1.03.0120, IP.1.1.03.0158 and IP.1.1.03.0039; and NextGenerationEU University grants: uniri-iz-25-6, uniri-iz-25-220, IIP\_010144, IIP\_010136}

\dataavailability{The AccessTrap-CS corpus and evaluation scripts are available at [repository URL to be added upon acceptance].}

\conflictsofinterest{The authors declare no conflicts of interest.}

%%%%%%%%%%%%%%%%%%%%%%%%%%%%%%%%%%%%%%%%%%
\abbreviations{Abbreviations}{
The following abbreviations are used in this manuscript:

\noindent
\begin{tabular}{@{}ll}
LLM   & Large Language Model \\
CS    & Computer Science \\
SMR   & Strategic Misconception Rate \\
DAR   & Dangerous Analogy Rate \\
OER   & Oversimplification Error Rate \\
ADHD  & Attention Deficit Hyperactivity Disorder \\
ANOVA & Analysis of Variance \\
HSD   & Honestly Significant Difference \\
HPC   & High Performance Computing \\
MoE   & Mixture of Experts \\
\end{tabular}
}

%%%%%%%%%%%%%%%%%%%%%%%%%%%%%%%%%%%%%%%%%%
\appendixtitles{no}
\appendixstart
\appendix

\section[\appendixname~\thesection]{CS Concepts and Subfields}
\label{appendix:concepts}

% TODO: Popuniti s popisom 30 koncepata i njihovih subpolja.

\section[\appendixname~\thesection]{Prompt Templates}
\label{appendix:prompts}

% TODO: Popuniti s kompletnim promptovima za tri uvjeta.

\section[\appendixname~\thesection]{Evaluation Codebook}
\label{appendix:codebook}

% TODO: Popuniti s detaljnim rubricom za LLM-as-judge.

\section[\appendixname~\thesection]{Annotated Examples of Dangerous Analogies}
\label{appendix:analogies}

% TODO: Popuniti s primjerima opasnih analogija i ekspertnim objašnjenjima.

%%%%%%%%%%%%%%%%%%%%%%%%%%%%%%%%%%%%%%%%%%
\begin{adjustwidth}{-\extralength}{0cm}

\reftitle{References}

\bibliography{references}

\PublishersNote{}
\end{adjustwidth}
\end{document}